\chapter{Conclusion}
\label{chap:conclusion}
Real processes change while event logs capture them. Such a concept drift is not only about the order of activities, but also about the resources and the situation of all running cases. Existing detectors, however, mainly focus on control-flow or follow single KPIs. Also, existing public test data contain neither resource drifts nor arrival rate drifts (Section~\ref{sec:related_work:positioning}). Therefore, this thesis investigates whether a process can be monitored as a combination of intra-case, resource and inter-case states and whether this combination is a useful indicator for a concept drift (Section~\ref{sec:intro_ssec:rq}).

Hence, we first separate the terms drift type (what changes), drift perspective (which perspective a drift type belongs to) and drift mode (how the change affects the process over time). Our state-based framework adapts the execution states of~\textcite{berti2025ocel} from objects to cases and extends them to resource and inter-case states of feature windows. Each perspective passes separately through four steps: feature extraction, scaling and PCA, clustering into states with a SOM, k-means or DBSCAN as well as drift signals (Chapter~\ref{chap:approach}). The drift signal compares the state distribution of each signal window with a reference using a divergence. Since the drift signal is a time series, it is also suitable as input to change point detection algorithms like PELT \autocite{adams2023explainable} or ADWIN \autocite{baier2020handling}. Rheon generates synthetic event logs with nine drift types from all three perspectives, each sudden or gradual, and lists the injected drifts in a metadata file (Chapter~\ref{chap:rheon}). Our library kairo implements the framework, and in our dashboard, users interpret the results with the LLM-based Kairo Copilot (Chapter~\ref{chap:impl}). Thus, we achieve~\ref{item:intro:rg_terminology} to~\ref{item:intro:rg_dashboard}.

For the research goals~\ref{item:intro:rg_evaluation} and~\ref{item:intro:rg_limitations}, we evaluate our framework on the 34 logs of the Padella benchmark and on three sets of 15 Rheon logs each (Chapter~\ref{chap:eval}). The Padella benchmark contains 29 control-flow and duration drifts in five real base logs, whereas the Rheon logs contain 103 drifts from all three perspectives. On each set, we compare two runs: in Run~1, agents created the configurations and thresholds, and in Run~2, we tuned them by hand. The results answer our research questions as follows (Section~\ref{sec:discussion:rq}):
\begin{itemize}
	\item \ref{item:intro:rq_states}~\emph{States:} The states show what changed. E.g., in the Padella log Bank account closure, each intra-case state of Run~1 corresponds to one activity, so that the deletion of an activity is directly visible.
	\item \ref{item:intro:rq_detection}~\emph{Detection:} With the configurations of Run~2, our framework detects 28 of the 29 Padella drifts and 93 of the 103 Rheon drifts. Its F1-scores lie between 0.862 and 0.933 on all four sets (Figure~\ref{fig:discussion:f1}).
	\item \ref{item:intro:rq_characterization}~\emph{Characterization:} Since each perspective has its own states, simultaneous drifts of different perspectives do not hide each other. However, some drift types also reach the features of other perspectives. Thus, the perspective of a drift only follows from comparing all three drift signals. A sudden drift shows as a spike against the previous signal window, whereas a gradual drift shows better in longer signal windows with a reference over several signal windows (Section~\ref{sec:eval:padella:gradual}). Section~\ref{sec:approach:features} shows to which drift types the features of each perspective respond, e.g., time drifts need the durations per resource (Section~\ref{sec:eval:padella:time_fast}).
	\item \ref{item:intro:rq_applicability}~\emph{Applicability:} On the same logs, the configurations of the agents only reach F1-scores from 0.309 to 0.774. A drift that no chosen feature captures stays hidden. Moreover, real logs contain changes that a benchmark does not know, e.g., at the turn of the year 2017/18 in two Padella base logs.
\end{itemize}

Therefore, we answer our main research question with yes. A process can be monitored as a combination of intra-case, resource and inter-case states, and their changes are a useful and explainable indicator for concept drift. However, this only holds if the chosen features capture the drift and the configuration fits the log. Our agent skill and the Copilot help to find such configurations, but do not replace a human (Section~\ref{sec:discussion:llm}).

\paragraph{Future Work.}
Based on our results, we see five directions for future work:
\begin{itemize}
	\item \emph{Automatic configuration:} Since the configuration determines the results the most, choosing it automatically is the most important next step. The rules of thumb in Section~\ref{sec:discussion:config} are a starting point for that.
	\item \emph{Online drift detection:} On an event stream, scaling, PCA and clustering learn their models on a warm-up period. Afterwards, the framework only assigns each new row to a state. Without the mean of all signal windows as reference (Section~\ref{sec:approach:drift_signals}) and without features that know the future of a case, such as the case progress, our dashboard then reports drifts live. After a drift, e.g., with new activities, the framework learns the models again.
	\item \emph{Object-centric drift detection:} \textcite{berti2025ocel} compute execution states for objects in object-centric event logs (OCELs). Applied to OCELs, our framework computes states and a drift signal per object type. \textcite{adams2023explainable} already detect drifts in OCELs, but in single time series instead of states. For the evaluation, we plan to extend Rheon to object-centric logs with drifts (Section~\ref{sec:rheon:limitations}).
	\item \emph{Rheon as training data:} \textcite{kraus2025machine} train the supervised detector CV4CDD-4D on 40,000 logs with control-flow drifts generated by CDLG and plan to include the time, resource and data perspectives. Rheon logs provide such training data from all three perspectives, and their metadata files serve as labels.
	\item \emph{State distances:} The divergences only compare the shares of the states (Section~\ref{sec:approach:drift_signals}). Thus, a transition into a similar state counts as much as a transition into a very different state. As for the trace distributions of \textcite{brockhoff2020time}, the Earth Mover's Distance additionally weighs each shift by the distance between the states.
\end{itemize}
